\documentclass[conference]{IEEEtran}
\IEEEoverridecommandlockouts
\usepackage{cite}
\usepackage{amsmath,amssymb,amsfonts}
\usepackage{algorithmic}
\usepackage{graphicx}
\usepackage{textcomp}
\usepackage{booktabs}
\usepackage{xcolor}
\usepackage{hyperref}
\usepackage{float}

\def\BibTeX{{\rm B\kern-.05em{\sc i\kern-.025em b}\kern-.08em
    T\kern-.1667em\lower.7ex\hbox{E}\kern-.125emX}}
\begin{document}


\title{Near-Real-Time ISAC-Based Jammer Detection and Localization using OCUDU\\}

\author{
    % Group all 4 authors in a single block
    \IEEEauthorblockN{
        Samanvitha Balijepalli\IEEEauthorrefmark{1}, 
        Aniket Jadhav\IEEEauthorrefmark{1}, 
        Sandra Mansour\IEEEauthorrefmark{1}, and 
        Rahman Doost-Mohammady\IEEEauthorrefmark{2}
    }
    
    % Provide the shared department and university details
    \IEEEauthorblockA{
        Department of Electrical and Computer Engineering \\
        Rice University, Houston, USA \\
        % Group or list the emails cleanly linked by the refmarks
        Email: \IEEEauthorrefmark{1}\{sb273, aj156, sm367\}@rice.edu, 
               \IEEEauthorrefmark{2}doost@rice.edu
    }
    }

\maketitle

\begin{abstract}
Malicious jammers threaten wireless communication systems, yet detecting and
localizing them often requires dedicated RF sensing hardware that is costly to
deploy and difficult to scale. This work presents a software-defined 5G
framework for jammer detection and localization using OCUDU together with
existing communication infrastructure. The first phase introduced low-overhead
sensing opportunities within the 5G NR frame and demonstrated reliable jammer
detection, with jammer-power-based detection achieving an F1-score of $1.0$ and
SNR-based detection achieving $92.18\%$ accuracy. The system has since been
extended to a live 5G NR testbed operating on band n78 using a USRP~B210, a
Quectel 5G modem, and an ADALM-Pluto jammer. Current localization work focuses
on multi-antenna direction estimation using ESPRIT, with MUSIC retained as a
baseline. The proposed architecture uses two spatially separated four-antenna
sensing stations to estimate jammer azimuth and elevation and combine these
measurements to determine the jammer's three-dimensional position $(x, y, z)$.
\end{abstract}

\begin{IEEEkeywords}
Integrated Sensing and Communication (ISAC), jammer detection, jammer
localization, 5G NR, wireless sensing, threshold-based detection,
ESPRIT, angle-of-arrival
\end{IEEEkeywords}


% -----------------------------------------------------------------------
\section{Introduction}
% -----------------------------------------------------------------------

Malicious radio-frequency (RF) jammers pose an increasingly significant threat to modern wireless systems. By transmitting continuous-wave, broadband, or other forms of interference, jammers can degrade LTE/5G communication links and disrupt safety-critical wireless services. Conventional jammer detection and localization systems often rely on dedicated spectrum-monitoring receivers, direction-finding antennas, and other specialized RF hardware. Although effective, these approaches increase deployment cost and can be difficult to scale or adapt as communication systems evolve. This motivates the development of software-defined sensing approaches that reuse existing communication infrastructure for jammer detection and localization.

Integrated Sensing and Communication (ISAC) provides a framework in which communication infrastructure can simultaneously support data transmission and environmental sensing. By reusing existing radio resources and hardware, ISAC can reduce the need for separate sensing infrastructure while enabling functions such as interference detection and source localization. Open and programmable 5G platforms are particularly useful for this purpose because they provide access to radio scheduling and physical-layer measurements that are typically hidden in commercial systems. The current implementation uses OCUDU~\cite{ocudu_project} together with Open5GS~\cite{open5gs_site} to provide a programmable 5G environment for jammer sensing and localization.

The first phase of this work developed a low-overhead jammer sensing and detection framework by introducing periodic quiet sensing opportunities within the 5G NR frame~\cite{3gpp38211}. During these intervals, the receiver can observe interference without expected user-equipment (UE) transmission. The resulting framework demonstrated reliable jammer detection, with jammer-power-based detection achieving an F1-score of 1.0 and SNR-based detection achieving $92.18\%$ accuracy. The system has since been extended from a software-based environment to a live 5G NR testbed operating on band n78 using a USRP B210, a Quectel 5G modem as the UE, and an ADALM-Pluto as the jammer source.

The current phase focuses on extending jammer sensing toward three-dimensional localization. Complex multi-antenna RF measurements are collected to estimate the direction from which the jammer signal arrives. MUSIC is retained as a baseline direction-of-arrival (DoA) algorithm, while ESPRIT is being investigated as the primary approach because it estimates direction using the rotational relationship between spatially shifted antenna subarrays without requiring an angular grid search. The proposed architecture uses two spatially separated sensing stations, each ultimately equipped with four receive antennas, to estimate jammer azimuth and elevation. The directional measurements from the two stations are then combined to estimate the jammer's three-dimensional position $(x,y,z)$.

The main contributions of this work are: (i) a low-overhead, scheduler-based sensing framework for jammer detection within a 5G NR system; (ii) validation of jammer sensing and AoA behavior on a live 5G NR hardware testbed; (iii) investigation of ESPRIT-based multi-antenna direction estimation, with MUSIC retained as a reference method; and (iv) development of a two-station architecture for three-dimensional jammer localization. Together, these components provide a path toward near-real-time jammer detection and localization using programmable 5G infrastructure with limited additional sensing hardware.


\subsection*{Related Work}
Chiarello et al.~\cite{9569600} proposed a jammer-detection technique based on
subcarrier blanking for 5G and beyond systems. In their approach, selected
time-frequency resources are intentionally left unused by legitimate
transmissions, allowing unexpected energy received in the blanked resources to
indicate jammer activity. Their work demonstrates the usefulness of creating
silent communication resources for interference sensing, but focuses on jammer
detection rather than determining the jammer's physical location.

For localization, Menta et al.~\cite{article1} investigated AoA-based
localization in 5G ultra-dense networks using multiple Transmission-Reception
Points (TRPs). Each TRP estimates the AoA of an uplink signal, and the resulting
measurements are combined using an Extended Kalman Filter to estimate and track
the transmitter's two-dimensional position. This work demonstrates the
feasibility of combining directional measurements from multiple cellular
sensing locations, although it considers a cooperative UE rather than an
uncooperative jammer and focuses primarily on two-dimensional localization.

Estimating a jammer's complete spatial direction also requires both azimuth and
elevation information. Dong et al.~\cite{article2} proposed a two-dimensional
DoA estimation approach using an L-shaped antenna array consisting of two
perpendicular uniform linear arrays. The geometry enables estimation of both
azimuth and elevation from the received signals. Their work demonstrates the
use of multi-dimensional antenna arrays for obtaining the directional
information needed for three-dimensional localization, although it focuses on
DoA estimation rather than determining the complete $(x,y,z)$ position of a
source using multiple sensing stations.

High-resolution subspace methods provide several approaches for obtaining these
directional estimates. MUSIC, introduced by Schmidt~\cite{1143830}, separates
the signal and noise subspaces of the array covariance matrix and evaluates a
pseudospectrum to determine likely source directions. Root-MUSIC~\cite{45540}
replaces the angular grid search with polynomial rooting, while
ESPRIT~\cite{32276} exploits rotational invariance between shifted antenna
subarrays to estimate direction directly. The reduced need for angular
searching makes ESPRIT particularly attractive for near-real-time localization,
although its effectiveness depends on suitable antenna geometry and accurate
phase relationships between receive channels. Sparse reconstruction
techniques~\cite{inbook} provide another approach to DoA estimation but can
introduce greater computational complexity.

SpotFi~\cite{article} demonstrated that complex OFDM measurements across
antennas and subcarriers can be used for localization and multipath resolution.
SpotFi exploits both spatial and frequency-domain phase information to
distinguish propagation paths and identify a likely direct path. However,
SpotFi operates in a cooperative Wi-Fi environment where channel state
information is available from normal communication transmissions. In the
jammer-localization problem considered here, the jammer is uncooperative and
conventional UE channel estimates may not be available during quiet sensing
resources. The proposed system therefore operates on complex multi-antenna RF or
frequency-domain resource-grid measurements that preserve the relative phase
information required for DoA estimation.

The proposed work combines concepts from these prior approaches. Quiet
communication resources provide opportunities for jammer sensing, multi-antenna
processing provides azimuth and elevation estimates, and measurements from two
spatially separated sensing stations are combined to determine the jammer's
three-dimensional position. Unlike prior work that addresses only jammer
detection, cooperative UE localization, or single-array DoA estimation, the
objective is to integrate these components into a near-real-time 5G jammer
detection and three-dimensional localization framework.

\begin{figure}[H]
    \centering
    \includegraphics[width=0.48\textwidth]{Abstract.png}
    \caption{Joint communication and sensing framework for jammer detection
    and localization using data-scheduling-based sensing opportunities in a
    5G gNB. Phase 1 established jammer detection. The current phase extends the framework to multi-antenna direction estimation and two-station 3D localization using ESPRIT, with MUSIC retained as a baseline.}
    \label{fig:system_framework}
\end{figure}


% -----------------------------------------------------------------------
\section{Methodology}
% -----------------------------------------------------------------------

\subsection{System Architecture}

The proposed system is built on the OCUDU~\cite{ocudu_project} 5G software
stack and uses communication-driven sensing opportunities to detect and
eventually localize jammer activity without dedicated sensing hardware.
The full platform consists of an Open5GS~\cite{open5gs_site} core network,
a modified OCUDU gNB running on a USRP~B210 radio, a Quectel 5G modem acting
as the UE, and an ADALM-Pluto jammer. All signals, including uplink traffic
and jammer interference, are exchanged over the air, as shown in
Fig.~\ref{fig:system_flowchart}.

The gNB was modified to create sensing opportunities, described in the next
subsection, during which the UE is silent and the receiver captures ambient RF
energy, interference, and jammer-induced signal components. The
frequency-domain samples from these slots are stored for post-processing. Each
recording is divided into jammer-OFF and jammer-ON phases according to the time
at which the jammer is switched on, providing the ground-truth labels needed
for detection and localization evaluation.

A post-processing layer converts the captured grids and logs into data
structures for offline analysis. In the current stage, sensing is performed
using the interference floor and the spatial structure of the received signal
across the two receive antennas, processed by lightweight threshold-based
detectors and a phase-difference angle-of-arrival (AoA) estimator. The
longer-term goal is to extend this pipeline toward four-antenna ESPRIT
direction estimation and three-dimensional localization.

\begin{figure}[htbp]
    \centering
    \includegraphics[width=0.8\linewidth]{System Architecture.png}
    \caption{System architecture and processing flow for the proposed
    sensing framework.}
    \label{fig:system_flowchart}
\end{figure}

\subsection{Sensing-Slot Capture}

To observe the jammer without dedicated sensing hardware, we need uplink
samples in which the UE is silent, so that only noise, interference, and the
jammer are present. We used two mechanisms in the gNB.

\textbf{Scheduler-reserved symbols.} We modified the gNB scheduler so that, on
every fourth slot, uplink grants end at symbol~10, leaving the remaining
symbols free for sensing. The reservation is enabled by an environment flag and
was confirmed to fire in the same slot of each frame. In our captures, however,
the reserved symbols did not become UE-silent: with the single-layer (rank-1)
uplink of our UE, the scheduler repacked the data into fewer symbols instead of
leaving a gap, so no jammer-only symbols appeared.

\textbf{Per-slot uplink tap.} We therefore added a read-only tap inside the
gNB's upper PHY that fires on every received uplink slot, scheduled or not, and
records the full frequency-domain grid (612 subcarriers on both receive ports)
once per slot. Slots in which the UE has nothing to send contain only noise,
interference, and the jammer. Spacing the uplink traffic (e.g.\ pings at
0.3~s intervals) makes about 43\% of slots idle, so the live gNB itself
supplies jammer-only samples without any scheduler change.

This second mechanism depends on spare uplink capacity, which will not exist
under sustained traffic. Reliable sensing in production therefore requires the
gNB to reserve whole uplink slots for sensing, and extending the scheduler
reservation from symbols to slots is part of the planned work.

\subsection{Detection and Localization Processing Framework}
\label{sec:aoa}

Let $Y_{p,k,s}$ denote the received resource element on receive port
$p\in\{0,1\}$, subcarrier $k\in\{1,\dots,612\}$, and uplink slot $s$. The
energy of slot $s$ is
\begin{equation}
P_s=\frac{1}{2\cdot 612}\sum_{p=0}^{1}\sum_{k=1}^{612}\lvert Y_{p,k,s}\rvert^{2}.
\end{equation}
Let $\mathcal S_q$ be the set of the fraction $q$ of slots with the lowest
energy in a recording, which are the UE-silent slots. The interference floor
and its rise between the jammer-OFF and jammer-ON phases are
\begin{equation}
P_{\mathrm{floor}}=\frac{1}{\lvert\mathcal S_q\rvert}\sum_{s\in\mathcal S_q}P_s,
\qquad
\Delta P=10\log_{10}\frac{P_{\mathrm{floor}}^{\mathrm{ON}}}{P_{\mathrm{floor}}^{\mathrm{OFF}}}.
\end{equation}

Define the two-antenna snapshot $\mathbf x_{k,s}=[Y_{0,k,s},\,Y_{1,k,s}]^{T}$.
The spatial covariance matrix over the UE-silent slots is
\begin{equation}
\hat{\mathbf R}=\frac{1}{K}\sum_{s\in\mathcal S_q}\sum_{k}\mathbf x_{k,s}\mathbf x_{k,s}^{H},
\end{equation}
where $K$ is the number of snapshots. With eigenvalues
$\lambda_1\ge\lambda_2$ of $\hat{\mathbf R}$, we use the eigenvalue ratio and
the spatial coherence,
\begin{equation}
\rho=10\log_{10}\frac{\lambda_1}{\lambda_2},
\qquad
\gamma=\frac{\lvert\hat R_{12}\rvert}{\sqrt{\hat R_{11}\hat R_{22}}}.
\end{equation}
A large $\rho$ and a high $\gamma$ indicate a single spatially coherent source,
which in UE-silent slots is the jammer.

For a feature $f\in\{P_{\mathrm{floor}},\gamma\}$, the class statistics are
the sample mean and standard deviation of the OFF and ON recordings,
\begin{equation}
\mu_f=\frac{1}{N}\sum_{i=1}^{N}f_i,\qquad
\sigma_f=\sqrt{\frac{1}{N-1}\sum_{i=1}^{N}(f_i-\mu_f)^2}.
\end{equation}
Using the class means, the midpoint threshold and threshold detector are
\begin{equation}
T_f=\frac{\mu_{f,\mathrm{OFF}}+\mu_{f,\mathrm{ON}}}{2},\qquad
\hat y_f(i)=\begin{cases}1,& f_i>T_f\\0,&\text{otherwise,}\end{cases}
\end{equation}
and the combined detector declares a jammer if either feature detector fires,
$\hat y_{\mathrm C}(i)=\hat y_{P_{\mathrm{floor}}}(i)\lor\hat y_{\gamma}(i)$.

For direction estimation, the phase difference between the two receive chains
gives the jammer azimuth. With antenna spacing $d$ and wavelength $\lambda$,
\begin{equation}
\Delta\varphi=\arg\hat R_{12}-\varphi_0,\qquad
\theta=\arcsin\!\Big(\frac{\lambda\,\Delta\varphi}{2\pi d}\Big),
\end{equation}
where $\varphi_0$ is a fixed phase offset between the two receive chains caused
by cable and RF mismatch. We measure it once with the jammer at boresight,
$\varphi_0=\arg\hat R_{12}\big|_{\theta=0}$, and reuse it as long as the
cabling is unchanged.

Two antennas give a single bearing angle and coarse resolution. The next phase
uses two stations, each with a $2\times2$ planar array. For a source at polar
angle $\theta$ and azimuth $\phi$, let $u=\sin\theta\cos\phi$ and
$v=\sin\theta\sin\phi$. With wavenumber factor $\kappa=2\pi d/\lambda$, the
steering vector of element $(m,n)$ is
$a_{mn}=\exp\!\big(j\kappa(mu+nv)\big)$, $m,n\in\{0,1\}$. ESPRIT uses the
shift invariance between element pairs along each array axis: with signal
subspace $\mathbf E_s$ of $\hat{\mathbf R}$ and selection matrices
$\mathbf J_{x1},\mathbf J_{x2}$,
\begin{equation}
\boldsymbol\Psi_x=(\mathbf J_{x1}\mathbf E_s)^{\dagger}\mathbf J_{x2}\mathbf E_s,
\qquad u=\frac{\arg\lambda_x}{\kappa},
\end{equation}
where $\lambda_x$ is the eigenvalue of $\boldsymbol\Psi_x$, and likewise
$v=\arg\lambda_y/\kappa$ from $\boldsymbol\Psi_y$. This gives
$(\theta,\phi)$ in closed form without a grid search~\cite{32276}. Each
station $i$ at known position $\mathbf p_i$ then yields a unit direction vector
$\mathbf d_i$, and the jammer position is the least-squares point closest to
both rays,
\begin{equation}
\hat{\mathbf s}=\Big(\sum_{i=1}^{2}\big(\mathbf I-\mathbf d_i\mathbf d_i^{T}\big)\Big)^{-1}
\sum_{i=1}^{2}\big(\mathbf I-\mathbf d_i\mathbf d_i^{T}\big)\mathbf p_i,
\end{equation}
which gives $(x,y,z)$. The ESPRIT and triangulation stages are planned and have
not yet been evaluated.

\subsection{Test Environment}

All experiments were run indoors on a live 5G NR testbed in band n78 (center
frequency 3489.42~MHz, 20~MHz bandwidth, 30~kHz subcarrier spacing, TDD,
single-user scenario). The gNB is OCUDU running on a laptop with Linux OS,
connected to a USRP~B210 with two receive antennas, together with an Open5GS
core network. The UE is a Quectel RM520N-GL modem connected to a second
laptop, which generates uplink traffic with iperf3. The jammer is an
ADALM-Pluto transmitting broadband noise co-channel in the same band, with its
antenna co-polarized with the gNB antenna.

The jammer was placed about 30~cm from the gNB antennas and the UE about 2~m
from the gNB. To obtain an angle swing, the jammer was moved between two
positions while the gNB and UE stayed fixed. Runtime execution was monitored
through the gNB logs to verify gNB initialization, UE attachment, and jammer
activation.

\textit{Testbed limitation:} the B210 provides only two receive antennas, so
the current testbed yields a single bearing angle with coarse resolution, and
each test case was recorded once. The four-antenna radios planned for the next
phase remove the first limitation.

\textit{Planned hardware:} the next phase replaces the B210 sensors with two
radios, each with four antennas arranged as a $2\times2$ planar matrix, placed
at least about 5~m apart. These radios use an O-RAN 7.2 split, so raw IQ is
not available and the AoA pipeline will work from per-antenna frequency-domain
resource elements exposed on the fronthaul. As with the B210, a one-time
per-antenna phase calibration will be required.

% -----------------------------------------------------------------------
\section{Results}
% -----------------------------------------------------------------------

\subsection{Recorded Measurements}

Measurements were recorded on the hardware testbed described in
Section~II-D (Fig.~\ref{fig:testbed}), with the jammer about 30~cm from the gNB
and the UE about 2~m from the gNB. The main recording contains 9005 uplink
slots, of which about 43\% are idle. UE-silent slots were taken as the quietest
$q=20\%$ of slots in each recording ($q=30\%$ in the two-position analysis). Each test case was recorded once, so the
results below are single-run observations without repeats or confidence
intervals.

\begin{figure}[htbp]
    \centering
    \includegraphics[width=0.9\linewidth]{testbed_setup.png}
    \caption{Hardware testbed: USRP~B210 gNB radio, Quectel UE, and ADALM-Pluto
    jammer.}
    \label{fig:testbed}
\end{figure}

\subsection{Jammer Detection from Sensing Slots}

Table~\ref{tab:hw_detect} summarizes detection from UE-silent slots. In the
quietest 20\% of uplink slots, the interference floor rose from 0.0014 to
0.0053 when the jammer was switched on, a rise of 5.9~dB. In the jammer-ON recording, no slot fell as low as the jammer-OFF floor,
because the jammer lifted the floor of every slot. The spatial coherence between the two antennas rose from
approximately 0.12 to 0.9, and the covariance eigenvalue ratio reached up to
16.8~dB in the strongest run. An earlier raw-IQ capture, taken with the gNB
stopped, showed the same behavior: the eigenvalue ratio was 0.9~dB without the
jammer and 16.6--17.3~dB with the jammer at two positions. With one recording per condition, the threshold detectors defined in Section~II-C are not yet evaluated statistically, so the features are reported
directly.

\begin{table}[htbp]
\centering
\caption{Sensing-slot detection metrics, jammer OFF vs.\ ON (single run).}
\label{tab:hw_detect}
\begin{tabular}{l c c}
\hline
\textbf{Metric} & \textbf{OFF} & \textbf{ON} \\
\hline
Idle-slot interference floor & 0.0014 & 0.0053 \\
Floor rise & \multicolumn{2}{c}{5.9 dB} \\
Spatial coherence & $\approx$0.12 & $\approx$0.9 \\
Max.\ eigenvalue ratio & --- & 16.8 dB \\
\hline
\end{tabular}
\end{table}

\subsection{Jammer Localization from Phase Difference}

The hardware phase offset was calibrated once at $\varphi_0=-1.857$~rad
($-106.4^\circ$), with the jammer at boresight (coherence 0.953, eigenvalue
ratio 16.2~dB). The calibrated estimate was validated physically: a jammer
placed near boresight read approximately $7.8^\circ$.

The jammer was then moved between two positions. The calibrated AoA moved from
$-49.3^\circ$ to $+28.0^\circ$, a swing of about $77^\circ$ that follows the
physical move (Table~\ref{tab:hw_aoa}, Fig.~\ref{fig:hw_aoa}). Position~1 gives
a strong, clean measurement. At position~2 the jammer signal was weaker, which
lowers the coherence to 0.552 and the floor rise to 1.9~dB, so that estimate is
less reliable.

\begin{table}[htbp]
\centering
\caption{Calibrated AoA at two jammer positions (single run each).}
\label{tab:hw_aoa}
\begin{tabular}{l c c c}
\hline
\textbf{Position} & \textbf{AoA} & \textbf{Coherence} & \textbf{Floor rise} \\
\hline
1 & $-49.3^\circ$ & 0.788 & 4.5 dB \\
2 & $+28.0^\circ$ & 0.552 & 1.9 dB \\
\hline
\end{tabular}
\end{table}

\begin{figure}[htbp]
    \centering
    \includegraphics[width=0.9\linewidth]{aoa_positions.png}
    \caption{Calibrated AoA versus jammer position, showing the $77^\circ$ swing
    between the two positions.}
    \label{fig:hw_aoa}
\end{figure}
\section{Discussion}

\textbf{Dataset quality and measurement repeatability.} The structured hardware sweep provides labeled measurements for evaluating jammer detection and localization under controlled
over-the-air conditions. Residual variation may
arise from changes in multipath propagation, antenna placement, transmitter--receiver alignment, and hardware-level gain or phase differences. In particular, the indoor environment introduces
position-dependent reflections that can alter the observed signal power and spatial response even when the nominal configuration remains unchanged. These effects are important when interpreting detection performance and AoA estimates, since variations in the received signal may reflect both jammer activity and propagation conditions. Consequently, the dataset establishes an experimental baseline for assessing the robustness of software-defined sensing
under practical wireless conditions.


\textbf{MUSIC--ESPRIT trade-offs for AoA estimation.}
MUSIC and ESPRIT offer complementary approaches to estimating
the direction of an incident signal. MUSIC estimates AoA by
searching a spatial spectrum, allowing the angular search region
and resolution to be controlled explicitly. However, its
computational cost increases with the density of the angular
grid, and its performance can degrade under low-SNR conditions,
limited snapshots, and multipath-induced source correlation.
ESPRIT instead exploits rotational invariance between appropriately
constructed subarrays to estimate AoA without an explicit angular
grid search, making it attractive for computationally efficient
localization. With four array elements, the idealized MUSIC
formulation can accommodate up to three uncorrelated sources
under suitable assumptions, whereas ESPRIT using two-element
subarrays is limited to one identifiable source in the
corresponding single-source formulation. These limits depend on
the array model, signal-subspace rank, and estimation conditions.

 Furthermore, classical ESPRIT
requires an appropriate shift-invariant array structure, and
hardware calibration errors can violate the assumed array model.
The choice between the two estimators must therefore be informed
by both the measured performance and the antenna geometry used
in the deployment.

\textbf{Limitations and future work.}
The current evaluation is limited to an indoor environment and
a single jammer type, restricting the extent to which the results
can be generalized to different propagation conditions,
interference waveforms, and deployment scenarios. The
$2\times2$ antenna configuration also limits the available
spatial aperture and the ability to resolve closely spaced
directions or estimate elevation reliably. In addition,
two-dimensional position estimation from measurements at the
two receiving stations has not yet been fully evaluated.
Further work will investigate localization accuracy across
different jammer positions, signal-to-interference conditions,
and multipath configurations, while quantifying the trade-off
between estimation accuracy and computational cost for MUSIC
and ESPRIT. Integration of the detection and localization
pipeline into a Near-Real-Time RAN Intelligent Controller
(xApp) is ongoing. This integration will enable future
evaluation of end-to-end latency, sensing overhead, and
localization performance within an Open RAN-based ISAC
framework.

The entire source code and related images are available at \href{https://github.com/AniketttJadhav/srsRAN-Open5GS}{this GitHub repository}.

% -----------------------------------------------------------------------
\section*{Acknowledgments}
% -----------------------------------------------------------------------

This work was conducted as part of the Master of Electrical and Computer
Engineering program at Rice University. The authors acknowledge the
academic support and institutional resources provided by Rice University,
as well as the guidance and feedback from Dr.\ Rahman Doost throughout
the project. The statements made in this work are solely the
responsibility of the authors.



\section*{CRediT Author Statement}

\textbf{Sandra Mansour:} Software, Formal analysis, Visualization,
Writing - Original draft, Reviewing and Editing. \textbf{Samanvitha Balijepalli:}
Investigation, Algorithm analysis, Writing - Review \& Editing.
\textbf{Aniket Jadhav:} Investigation, Methodology, Software.
\textbf{Rahman Doost-Mohammady:} Conceptualization, Supervision, Project
administration, Resources.

\bibliographystyle{IEEEtran}
\bibliography{references}

\end{document}